\documentclass[11pt,a4paper]{article}
\usepackage[margin=1in]{geometry}
\usepackage{amsmath,amssymb}
\usepackage{enumitem}
\usepackage{hyperref}
\usepackage{booktabs}
\usepackage{parskip}

\title{Chess Thermodynamics:\\A Partition-Function Framework for Positional Evaluation}
\author{}
\date{}

\begin{document}
\maketitle

\begin{abstract}
We develop a chess evaluation framework in which thermodynamic quantities
arise from a genuine partition function over legal continuations.
Each chess position induces an ensemble of reachable futures. A Boltzmann
distribution over this ensemble, parameterized by a rationality
temperature~$T$, defines the partition function~$Z$, free energy~$F$,
continuation entropy~$S$, and specific heat~$C$. The leaf evaluator
is an \emph{ideal gas}: each piece contributes independently via its
x-ray mobility, weighted by an evolvable chemical potential. Positional
understanding---coordination, tactics, king safety---emerges from the
search tree rather than from fitted regression weights. Standard minimax
search is recovered as the zero-temperature limit $T \to 0$, while finite
temperature gives entropy genuine strategic value.
\end{abstract}

\section{Motivation}

The first version of this framework assigned thermodynamic names---$U$,
$T$, $S$, $F$, $PV$---to static board features. While implementable, this
approach was vulnerable to the criticism that the physics was decorative:
entropy was defined as x-ray redundancy rather than as a count of
accessible states, temperature was a derived statistic rather than a
Lagrange multiplier, and the ``potential energy'' $PV$ was a polynomial
regression with no independent pressure or volume.

The resolution is not to weaken the analogy but to make it structurally
exact. This means defining an ensemble, microstates, a partition function,
and deriving all thermodynamic quantities from it. The result is a
framework where the thermodynamic language does real work.

\section{Internal Energy: The Ideal Gas}

The leaf evaluator treats each piece as an independent particle in an
ideal gas on the $8 \times 8$ lattice. The internal energy of a side is
the weighted sum of x-ray mobilities:

\begin{equation}
U_c = \sum_{i \in \text{pieces of color } c} \mu_{t_i} \, f_i
\end{equation}

where $f_i$ is the x-ray mobility of piece~$i$ (the number of squares it
reaches on an otherwise empty board from its current square), $t_i$ is its
piece type, and $\mu_t$ is the \textbf{chemical potential} of type~$t$.

X-ray mobility is a function of piece type and square alone, independent
of the current board configuration. This makes $U$ a measure of the
\emph{intrinsic energy potential} of the pieces. With all chemical
potentials equal to~1, a queen on d4 contributes $\sim$27~squares and a
pawn on e2 contributes $\sim$4. Material advantage is directly reflected:
removing a piece always reduces $U$ for that side.

The leaf evaluation of a position is:

\begin{equation}
\mathcal{L}(P) = U_w(P) - U_b(P)
\end{equation}

from White's perspective. This is the ``ideal gas'' contribution---each
piece acts independently.

\subsection{Chemical Potentials}

In thermodynamics, the chemical potential $\mu_t = \partial F / \partial
N_t$ measures the marginal free energy cost of adding or removing one
particle of type~$t$. In our framework, the chemical potentials $\mu_t$
weight each piece type's contribution per unit of x-ray reach:

\begin{equation}
\mu_t \in \{\mu_P, \mu_N, \mu_B, \mu_R, \mu_Q, \mu_K\}
\end{equation}

With equal weights $\mu_t = 1$, the traditional piece-value ratios emerge
naturally from x-ray mobility: a knight on a central square has
$f \approx 8$, a rook has $f \approx 14$, a queen has $f \approx 25$. The
chemical potentials allow the framework to learn corrections---for example,
that a knight's per-square contribution is worth more than a queen's,
reflecting the knight's unique geometry.

These potentials are discovered by evolutionary self-play or fitted from
data, providing a dynamic, position-aware theory of material value.

\section{The Ensemble and Partition Function}

A chess position $P$ is not a thermodynamic state in isolation; it is a
\textbf{macrostate} that induces an ensemble of accessible futures. We
define the ensemble as the tree of legal continuations to depth~$h$:

\begin{equation}
\Omega_h(P) = \{\text{positions reachable from } P \text{ within } h
\text{ plies}\}
\end{equation}

Over this ensemble, we define a Boltzmann distribution parameterized by
the rationality temperature~$T$:

\begin{equation}
\pi_T(a \mid P) = \frac{\exp(Q(P, a) / T)}{Z_T(P)}
\end{equation}

where $a$ ranges over legal moves, $Q(P, a)$ is the value of the
resulting position (computed recursively), and $Z_T(P)$ is the partition
function:

\begin{equation}
Z_T(P) = \sum_{a \in \text{Legal}(P)} \exp(Q(P, a) / T)
\end{equation}

The free energy of the position is:

\begin{equation}
\boxed{F_T(P) = T \log Z_T(P) = T \log \sum_{a} \exp(Q(P, a) / T)}
\end{equation}

This is the \textbf{log-sum-exp} or \textbf{soft-max} operator. For a
two-player game with alternating optimization, we use the negamax
convention: $Q(P, a) = -F_T(P_a)$, where $P_a$ is the position after
move~$a$ and the sign flip accounts for the change of perspective.

\subsection{The Zero-Temperature Limit}

As $T \to 0$, the Boltzmann distribution concentrates on the
highest-value move, and the free energy converges to the minimax value:

\begin{equation}
\lim_{T \to 0} F_T(P) = \max_{a} Q(P, a) = \max_{a} \{-F_0(P_a)\}
\end{equation}

This is precisely the standard minimax algorithm. \textbf{Minimax is the
zero-temperature limit of thermodynamic evaluation.} Finite temperature
generalizes minimax by giving entropy---having many good options---genuine
strategic value.

\section{Continuation Entropy}

The Shannon entropy of the Boltzmann distribution over legal moves is:

\begin{equation}
S_T(P) = -\sum_{a} \pi_T(a \mid P) \log \pi_T(a \mid P)
\end{equation}

This is the \textbf{continuation entropy}: a measure of how many
strategically viable futures the position has, weighted by their quality.

Chess interpretations:

\begin{itemize}[nosep]
\item \textbf{High entropy}: Many viable continuations; flexible;
  strategically rich. The side has options.
\item \textbf{Low entropy}: Few viable continuations; forced; tactically
  constrained. The position is concrete.
\item \textbf{Entropy collapse}: A forcing line, mating net, or
  liquidation sequence. The number of playable moves drops sharply.
\item \textbf{Entropy advantage}: One side has more playable futures
  than the other---a thermodynamic formulation of the initiative.
\end{itemize}

\section{Temperature as Rationality}

Temperature $T$ is the Lagrange multiplier controlling the tradeoff
between energy optimization and option diversity:

\begin{itemize}[nosep]
\item \textbf{Low $T$} $\to$ the player minimizes energy aggressively,
  choosing the single best move. This is concrete, forcing, engine-like
  play with low tolerance for inferior moves.
\item \textbf{High $T$} $\to$ the player explores more states, weighting
  many moves similarly. This is strategic, flexible, diffuse play where
  many continuations are considered viable.
\end{itemize}

The free energy satisfies the fundamental relation:

\begin{equation}
F_T = \langle Q \rangle + T \, S
\end{equation}

where $\langle Q \rangle = \sum_a \pi_T(a) \, Q(P, a)$ is the expected
value. The entropy term $T S$ is the \emph{option premium}: the free
energy exceeds the expected value because having choices is itself
valuable. A position is thermodynamically favorable not only when the
best move is good, but when the side has many good moves.

\section{Specific Heat: Tactical Volatility}

The specific heat is the variance of continuation values divided by
$T^2$:

\begin{equation}
C_T(P) = \frac{\text{Var}_{\pi_T}(Q)}{T^2} = \frac{\sum_a \pi_T(a)
\left(Q(P,a) - \langle Q \rangle\right)^2}{T^2}
\end{equation}

Chess interpretations:

\begin{itemize}[nosep]
\item \textbf{High specific heat}: Tactically volatile position. Small
  perturbations in move choice cause large evaluation swings. The move
  values are widely dispersed.
\item \textbf{Low specific heat}: Strategically stable position. Many
  moves lead to similar outcomes. The position is robust.
\end{itemize}

A sacrifice, breakthrough, or mating attack can be understood as a
\textbf{phase transition}: the position's accessible state space
reorganizes abruptly, reflected in a spike in specific heat.

\section{Evaluation}

The engine evaluation is the free-energy difference between the two
sides, converted to pawn-equivalent units:

\begin{equation}
\boxed{\mathcal{E}(P) = \frac{F_T(P)}{\bar{f}_P}}
\end{equation}

where $F_T(P)$ is computed from White's perspective (positive favors
White) and $\bar{f}_P$ is the average pawn x-ray mobility ($\approx 3.5$
squares), providing a natural scale for piece-equivalent units.

The evaluation bar maps $\mathcal{E}$ to win probability via the logistic
sigmoid $\text{Win\%} = 1/(1 + e^{-\mathcal{E}})$, with no arbitrary
scaling constants.

For actual move selection, the engine uses alpha-beta search (the $T = 0$
limit) with iterative deepening, null-move pruning, and late move
reductions. The thermodynamic quantities ($F$, $S$, $C$) are computed
separately via soft-minimax at moderate depth for display and analysis.

\section{Tactics as Entropy Collapse}

A tactical sequence is not merely a high-energy event---it is often an
\emph{entropy collapse}. Before the tactic, many legal-looking
continuations exist. After the forcing move, only one or two survive.

\begin{equation}
\Lambda(P) = S(P) - S(P_{a^*})
\end{equation}

where $a^*$ is the best move. The \emph{forcingness} $\Lambda$ measures
the entropy reduction imposed on the opponent:

\begin{itemize}[nosep]
\item \textbf{Check} = local entropy collapse (one legal king move often
  dominates)
\item \textbf{Double attack} = entropy bottleneck (opponent cannot defend
  both threats)
\item \textbf{Pin} = constrained phase space (piece's legal moves
  restricted)
\item \textbf{Mating net} = total entropy collapse (all moves lead to
  mate)
\item \textbf{Zugzwang} = entropy exists but all states are bad (legal
  moves exist, but every one worsens the position)
\end{itemize}

\section{Future Directions}

\subsection{Virial Expansion}

The current leaf evaluator is a pure ideal gas: pieces contribute
independently. A natural extension is a \emph{virial expansion} adding
pairwise interaction terms:

\begin{equation}
U(P) = \sum_i \mu_{t_i} f_i + \sum_{i < j} B(t_i, t_j, r_{ij}) \, q_i
q_j + \cdots
\end{equation}

where $B$ is a second virial coefficient depending on piece types and the
relationship $r_{ij}$ (attacks, defends, pins, shares file, near king,
etc.), and $q_i$ is the piece's ``charge'' (x-ray mobility). This would
capture coordination, obstruction, and tactical coupling at the leaf level,
reducing the search depth needed for positional understanding.

\subsection{Pressure and Volume}

Define volume as the accessible phase space of a side:
$V_c = |\text{legally reachable squares for side } c|$.
Then pressure $P_c = -\partial F / \partial V_c$ measures how much the
free energy worsens when space is compressed. Pawn chains, blockades, and
space advantages become thermodynamic quantities.

\subsection{Phase Detection}

Phase transitions (opening $\to$ middlegame, quiet $\to$ tactical) can be
detected from discontinuities in specific heat. The framework provides
intrinsic phase labels without external definitions.

\section{Summary}

\begin{center}
\renewcommand{\arraystretch}{1.3}
\begin{tabular}{@{}lll@{}}
\toprule
\textbf{Quantity} & \textbf{Definition} & \textbf{Chess meaning} \\
\midrule
$U$ & $\sum_i \mu_{t_i} f_i$ & Weighted x-ray mobility (ideal gas) \\
$T$ & Rationality parameter & Concrete ($T \to 0$) vs.\ flexible \\
$Z$ & $\sum_a \exp(Q_a / T)$ & Partition function over moves \\
$F$ & $T \log Z$ & Free energy (position value) \\
$S$ & $-\sum_a \pi_a \log \pi_a$ & Continuation entropy (optionality) \\
$C$ & $\text{Var}(Q) / T^2$ & Specific heat (tactical volatility) \\
$\mu_t$ & $\partial F / \partial N_t$ & Chemical potential (piece value) \\
$\mathcal{E}$ & $F / \bar{f}_P$ & Evaluation in pawn units \\
\bottomrule
\end{tabular}
\end{center}

The framework provides a structurally exact statistical-mechanical
evaluator. The partition function is real. The free energy is real. The
entropy is real. Minimax appears as the zero-temperature limit, and finite
temperature gives the initiative, flexibility, and optionality genuine
thermodynamic meaning.

\end{document}
